\subsection{Exceptional bandwidths in the equal-mass case}
\label{sec:exceptional-bandwidths}

For any fixed equal-mass site set that fails the reflection condition
of Corollary~\ref{cor:reflection}, only finitely many positive bandwidths
can preserve all strict decisions. Indeed, for some projection group
on some facet the balance in \eqref{eq:fiber} is a nonzero exponential
sum. Group by squared normal offset and discard zero coefficients to
write it as $\sum_{j=1}^m c_j e^{-s_j z}$, with $s_1<\cdots<s_m$ and
$z=1/(2\beta)$. It has at most $m-1$ distinct zeros for $z>0$. For the
induction, divide by $e^{-s_1 z}$ and differentiate. The derivative is
an exponential sum with $m-1$ nonzero terms, and Rolle's theorem bounds
the number of zeros of the divided sum by one plus the derivative's
count. A one-term sum has no zeros. Every preserving bandwidth must
be a zero of the chosen balance, regardless of the sampling horizon.

\paragraph{A rhombus preserving decisions at a single bandwidth.}
Take the four equal-mass sites $0,p,q,p+q$, where $p=(2,0)$ and
$q=(1/2,\sqrt{15}/2)$. The neighbors are the four sides and the short
diagonal $[p,q]$. A hull side gives a Voronoi facet because a point
sufficiently far along its outward normal from the side midpoint is
equidistant from its two endpoints and farther from the other sites.
The midpoint $(p+q)/2$ is nearer $p,q$ than $0,p+q$ because
$p^\top q>0$. The long diagonal is inactive: along the bisector of
$0,p+q$, at least one of $p,q$ is closer, since the difference between
the two pairs' summed squared distances is $-2p^\top q$.

On the side bisector $x_1=1$, the pair $0,p$ cancels in its projection
group. The pair $q,p+q$ has offsets $-1/2,3/2$. Its balance vanishes
if and only if
\[
 e^{-1/\beta}=\tfrac13,
 \qquad\text{that is, }\quad \beta=1/\log 3.
\]
Reflection across the line through $0,p+q$ exchanges $p,q$, and
central inversion exchanges opposite sides. Thus all four side
balances hold at the same bandwidth. The short-diagonal bisector is
the line through $0,p+q$, so it is balanced by reflection at all
bandwidths. Theorem~\ref{thm:flat} gives preservation at all horizons
at $\beta=1/\log 3<D^2/4=5/2$. Reflection in $x_1=1$ does not permute
the upper pair, so the orbit condition fails. The same side balance
excludes preservation at any other bandwidth.
